% !TEX program = pdflatex
% Presentation superviseur -- Week 1 -- 2026-04-13
% Knowledge Graph juridique francais

\documentclass[aspectratio=169,10pt]{beamer}

\usepackage[utf8]{inputenc}
\usepackage[T1]{fontenc}
\usepackage[french]{babel}
\usepackage{lmodern}
\usepackage{booktabs}
\usepackage{array}
\usepackage{tikz}
\usepackage{xcolor}
\usetikzlibrary{positioning, arrows.meta, shapes.geometric, fit, backgrounds}

\usetheme{Madrid}
\usecolortheme{seahorse}
\setbeamertemplate{navigation symbols}{}
\setbeamertemplate{footline}[frame number]
\setbeamerfont{footnote}{size=\tiny}

% Couleurs custom
\definecolor{accent}{RGB}{41, 92, 155}
\definecolor{accentsoft}{RGB}{230, 238, 247}
\definecolor{warn}{RGB}{200, 80, 40}

% Boxes
\newcommand{\keybox}[1]{%
  \begin{center}
  \colorbox{accentsoft}{\parbox{0.9\linewidth}{\centering #1}}
  \end{center}}

% Metadata
\title[KG juridique FR]{Construction d'un Knowledge Graph juridique français}
\subtitle{Point d'avancement --- Week 1}
\author{Matthieu Kaeppelin}
\institute{Stage FE Recherche}
\date{13 avril 2026}

\begin{document}

% ============================================================
\begin{frame}
  \titlepage
\end{frame}

% ============================================================
\begin{frame}{Plan}
  \tableofcontents
\end{frame}

% ============================================================
\section{Rappel du besoin}

\begin{frame}{Comment raisonne un juriste ? (1/2)}
  \keybox{Pour gagner un argument, il faut \textbf{3 choses} : une \textbf{règle}, des \textbf{faits}, et une \textbf{démonstration} que la règle s'applique aux faits. La \textbf{jurisprudence} fait le pont.}

  \vspace{0.3em}
  \textbf{Exemple --- le voisin bruyant}
  \begin{columns}[T]
    \column{0.55\textwidth}
    \footnotesize
    \begin{itemize}
      \item \textit{Faits} : mon voisin fait du bruit tous les soirs à 23h
      \item \textit{Règle} (art. 1240 C. civ.) : "toute faute causant un dommage oblige à réparer"
      \item \textit{Problème} : est-ce que "bruit à 23h" = une faute ?
      \item \textit{Jurisprudence} : Cass. 2018 --- les nuisances sonores répétées après 22h constituent un trouble anormal
      \item $\Rightarrow$ \textit{Argument} : règle + JP + faits = conclusion
    \end{itemize}

    \column{0.45\textwidth}
    \centering
    \footnotesize\textbf{Analogie}

    \vspace{0.3em}
    \renewcommand{\arraystretch}{1.15}
    \begin{tabular}{l l}
      \toprule
      \scriptsize\textbf{Science} & \scriptsize\textbf{Droit} \\
      \midrule
      \scriptsize Théorie            & \scriptsize Règle (article) \\
      \scriptsize Expériences passées & \scriptsize Arrêts passés (JP) \\
      \scriptsize Observation         & \scriptsize Faits de l'espèce \\
      \scriptsize Prédiction validée  & \scriptsize Argument juridique \\
      \bottomrule
    \end{tabular}
  \end{columns}

  \vspace{0.4em}
  \begin{exampleblock}{Important}
    \small Un juriste qui plaide = un chercheur qui doit citer \textbf{tous les articles et toutes les expériences passées pertinentes} pour convaincre.
  \end{exampleblock}
\end{frame}

\begin{frame}{Architecture du raisonnement \& outil cible (2/2)}
  \begin{columns}[T]
    \column{0.5\textwidth}
    \textbf{Architecture réelle (un graphe)}
    \vspace{0.2em}

    \centering
    \begin{tikzpicture}[scale=0.75, transform shape,
        arg/.style={rectangle, rounded corners, draw=accent, fill=accentsoft, thick, minimum width=2cm, minimum height=0.7cm, font=\small},
        art/.style={rectangle, draw=accent!80, thick, minimum width=1.6cm, minimum height=0.6cm, font=\scriptsize},
        jp/.style={ellipse, draw=accent!80, thick, minimum width=1.6cm, minimum height=0.6cm, font=\scriptsize}]
      \node[arg] (a) at (0,2.3) {Argument};
      \node[art] (art1) at (-2.2,0.6) {Article};
      \node[art] (art2) at (2.2,0.6) {Article};
      \node[jp] (jp1) at (-2.2,-1.1) {Arrêt};
      \node[jp] (jp2) at (2.2,-1.1) {Arrêt};
      \draw[-{Stealth}, thick] (a) -- (art1);
      \draw[-{Stealth}, thick] (a) -- (art2);
      \draw[-{Stealth}, thick] (a) -- (jp1);
      \draw[-{Stealth}, thick] (a) -- (jp2);
      \draw[-{Stealth}, dashed] (jp1) -- (art1) node[midway, left, font=\tiny] {applique};
      \draw[-{Stealth}, dashed] (jp2) -- (art2) node[midway, right, font=\tiny] {applique};
      \draw[<->, dashed] (jp1) -- (jp2) node[midway, below, font=\tiny] {confirme/casse};
    \end{tikzpicture}

    \vspace{0.3em}
    \footnotesize\textit{La JP ne fait pas que valider : elle interprète, qualifie les faits, et peut même créer la règle.}

    \column{0.5\textwidth}
    \textbf{Pourquoi c'est dur à automatiser}
    \begin{itemize}\footnotesize
      \item Un article seul ne suffit jamais
      \item 10 arrêts disent oui, 3 disent non $\rightarrow$ tout connaître
      \item Les arrêts se citent entre eux $\rightarrow$ un vrai réseau
      \item Des heures de travail manuel pour un juriste
    \end{itemize}

    \vspace{0.3em}
    \textbf{Outil cible}
    \keybox{\small Pour un argument donné, retrouver \textbf{les fondements} pertinents et \textbf{la jurisprudence} applicable (y compris contraire).}

    \textbf{Utilisateurs} : avocats, magistrats, juristes d'entreprise, particuliers\\
    \textbf{Gain} : temps $\cdot$ exhaustivité $\cdot$ traçabilité
  \end{columns}
  \vspace{1em}                                                                                        
  \footnotesize\textcolor{gray}{\textit{Vocabulaire : un \textbf{fondement} = un article de loi (ou un principe) $\cdot$ un \textbf{arrêt} = une décision de justice, composante de la jurisprudence.}}
\end{frame}

% ============================================================
\section{Problématiques}

\begin{frame}{Problématiques}
  \begin{columns}[T]
    \column{0.5\textwidth}
    \textbf{Problématique métier}
    \begin{itemize}\footnotesize
      \item Recherche juridique = raisonnement \textbf{multi-saut} (article $\rightarrow$ arrêt $\rightarrow$ article cité $\rightarrow$ JP interprétative)
      \item Le RAG vectoriel classique rate ces chaînes et ne restitue pas le raisonnement
      \item {Traçabilité} - sources citables
      \item {Transparence} - pourquoi cette remontée ?
    \end{itemize}

    \column{0.5\textwidth}
    \textbf{Problématique technique}
    \begin{itemize}\footnotesize
      \item Construire un KG juridique FR depuis Judilibre
      \item Garantir la qualité d'extraction LLM (Pas d'hallucinations)
      \item \textbf{Maintenabilité} : ajouter de nouveaux arrêts / articles sans tout reconstruire 
      \item \textbf{Temporalité} : evolution dans le temps des arrêts et articles
    \end{itemize}
  \end{columns}

  \vspace{0.5em}
  \begin{center}
    \textbf{Question de recherche}
  \end{center}
  \keybox{Dans quelle mesure une architecture GraphRAG hybride peut-elle améliorer la recherche juridique française en termes de \textbf{pertinence des arrêts}, \textbf{performance} et \textbf{explicabilité} du raisonnement ?}
\end{frame}

% ============================================================
\section{Données disponibles}

\begin{frame}{Données disponibles --- sources et éléments extractibles}
  \begin{columns}[T]
    \column{0.48\textwidth}
    \begin{block}{\small Articles de loi --- \textit{Légifrance}}
      \centering

      \vspace{0.4em}
      \footnotesize\textbf{Éléments identifiés}
      \begin{itemize}\scriptsize
        \item \textbf{Code} (Civil, Pénal, Travail\ldots)
        \item \textbf{Livre} / Titre / Chapitre / Section
        \item \textbf{Article} (numéro, libellé)
        \item \textbf{Contenu} (texte normatif)
        \item \textbf{Versions} (entrée en vigueur, abrogation)
      \end{itemize}
    \end{block}

    \column{0.48\textwidth}
    \begin{block}{\small Arrêts / Jurisprudence --- \textit{Judilibre}}
      \centering


      \vspace{0.4em}
      \footnotesize\textbf{Éléments identifiés}
      \begin{itemize}\scriptsize
        \item \textbf{Juridiction} (Cass., CE, CA\ldots) et formation
        \item \textbf{Date} et numéro de pourvoi
        \item \textbf{Articles cités} (fondements)
        \item \textbf{Moyens / arguments} des parties
        \item \textbf{Motifs} de la décision
        \item \textbf{Dispositif} (la décision elle-même)
      \end{itemize}
    \end{block}
  \end{columns}

  \vspace{0.3em}
  \keybox{\small Ces éléments structurés = \textbf{matière première} du KG juridique.}
\end{frame}

% ============================================================
\section{État de l'art}

\begin{frame}{État de l'art --- cartographie}
  \centering
  \begin{tikzpicture}[scale=0.9, transform shape,
      center/.style={circle, draw=accent, fill=accentsoft, thick, minimum width=2.2cm, align=center, font=\small\bfseries},
      branch/.style={rectangle, rounded corners, draw=accent!60, thick, minimum width=2.6cm, align=center, font=\footnotesize}]
    \node[center] (c) at (0,0) {KG juridique\\FR};
    \node[branch] (b1) at (-4.5,2) {KG juridiques\\{\tiny Belikov, D'Amato, LYNX}};
    \node[branch] (b2) at (0,2.5) {GraphRAG\\{\tiny Edge, Zhang, Peng}};
    \node[branch] (b3) at (4.5,2) {Construction KG + LLM\\{\tiny Bian 2025}};
    \node[branch] (b4) at (-4.5,-2) {NLP juridique FR\\{\tiny JuriBERT, CamemBERT}};
    \node[branch] (b5) at (0,-2.5) {Ontologies\\{\tiny ELI, AKN, LKIF, ECLI}};
    \node[branch] (b6) at (4.5,-2) {Données FR\\{\tiny Judilibre, Légifrance}};
    \foreach \x in {b1,b2,b3,b4,b5,b6} \draw[-, thick, accent!50] (c) -- (\x);
  \end{tikzpicture}

  \vspace{0.3em}
  \footnotesize \textbf{$\sim$50 sources cataloguées}, fichage ciblé en cours --- balayage large avant positionnement.
\end{frame}

% ============================================================
\section{GraphRAG --- concepts clés}

\begin{frame}{GraphRAG --- les 3 types de Knowledge Graph}
  \textit{D'après Zhang et al. 2025, Figure 2}

  \begin{columns}[T]
    \column{0.48\textwidth}
    \begin{block}{\small Indexation}
    \footnotesize
    DBpedia, YAGO, Wikidata\\[0.3em]
    KG généralistes servant d'\textbf{index} / référence transversale\\
    Larges, peu spécialisés\\
    Construction manuelle, encyclopédique
    \end{block}

    \column{0.48\textwidth}
    \begin{block}{\small Knowledge container}
    \footnotesize
    UMLS (médical), LYNX (droit EU)\\[0.3em]
    KG experts portant la \textbf{connaissance métier} dense\\
    Ontologie formelle\\
    Construction coûteuse (experts humains)
    \end{block}
  \end{columns}

  \vspace{0.6em}
  \begin{columns}
    \column{0.7\textwidth}
    \begin{alertblock}{\small Hybrid --- corpus-based \textit{(notre choix)}}
    \footnotesize
    Microsoft GraphRAG, LightRAG, HippoRAG\\[0.3em]
    KG \textbf{extrait par LLM} depuis un corpus de texte, éventuellement guidé par une ontologie légère\\
    $\Rightarrow$ souplesse du LLM $+$ structure d'une ontologie métier
    \end{alertblock}
  \end{columns}
\end{frame}

\begin{frame}{GraphRAG --- les 3 phases d'un pipeline}
  \textit{D'après Zhang et al. 2025, Figure 3}

  \vspace{0.5em}
  \begin{enumerate}
    \item \textbf{Knowledge Organization} (G-Indexing)
    \begin{itemize}\footnotesize
      \item chunking, extraction entités/relations, construction graphe, communautés
    \end{itemize}

    \item \textbf{Knowledge Retrieval} (G-Retrieval)
    \begin{itemize}\footnotesize
      \item matching requête $\rightarrow$ graphe
      \item stratégies : similarité sémantique, raisonnement logique, GNN, LLM, multi-round, hybride
    \end{itemize}

    \item \textbf{Knowledge Integration} (G-Generation)
    \begin{itemize}\footnotesize
      \item injection dans le LLM : in-context, fine-tuning (node/path/subgraph), refinement
    \end{itemize}
  \end{enumerate}

  \vspace{0.3em}
  \begin{exampleblock}{Positionnement des papiers pivots}
    \footnotesize
    \textbf{Edge 2024} : pipeline complet (1+3) $\bullet$ \textbf{Zhang 2025} : taxonomie (1+2+3)\\
    \textbf{Belikov-Raoult 2025} : phase 1 sur Cassation FR $\bullet$ \textbf{D'Amato 2025} : phase 1 + ontologie
  \end{exampleblock}
\end{frame}

% ============================================================
\section{Papiers lus en détail}

\begin{frame}{Papiers lus --- Microsoft GraphRAG}
  \begin{columns}
    \column{0.75\textwidth}
    \begin{block}{\small Edge et al. 2024 --- Microsoft GraphRAG}
    \footnotesize
    \textbf{Idée} : KG extrait par LLM + résumés de communautés Leiden $\rightarrow$ QFS map-reduce

    \vspace{0.3em}
    \textbf{Apports}
    \begin{itemize}\setlength\itemsep{-0.1em}
      \item Pipeline d'extraction complet (chunking $\rightarrow$ entités/relations $\rightarrow$ \textit{gleanings})
      \item Hiérarchie de communautés (Leiden récursif) + résumés multi-niveaux
      \item Inférence map-reduce pour sensemaking (QFS)
      \item Évaluation LLM-as-judge (comprehensiveness, diversity, empowerment)
    \end{itemize}
    \textbf{Limites}
    \begin{itemize}\setlength\itemsep{-0.1em}
      \item Pas de métrique de factualité / traçabilité
      \item Coût de construction élevé (multiples appels LLM)
    \end{itemize}
    \end{block}
  \end{columns}
\end{frame}

\begin{frame}{Papiers lus --- Surveys GraphRAG}
  \begin{columns}[T,onlytextwidth]
    \column{0.45\textwidth}
    \begin{block}{\small Zhang et al. 2025 --- Survey GraphRAG for Customized LLMs}
    \scriptsize
    \textbf{Idée} : taxonomie unifiée du GraphRAG en 3 phases pour LLMs spécialisés.

    \vspace{0.3em}
    \textbf{Apports}
    \begin{itemize}\setlength\itemsep{-0.1em}
      \item \textbf{3 types de KG} : open-domain, domain-specific, hybride
      \item \textbf{3 phases} : Knowledge Organization $\rightarrow$ Retrieval $\rightarrow$ Integration
      \item Table comparative des stratégies de retrieval (similarité, logique, GNN, LLM, multi-round)
      \item Cartographie des techniques d'intégration (in-context vs fine-tuning)
    \end{itemize}
    \textbf{Rôle} : \textbf{boussole} pour classer toutes les approches
    \end{block}

    \column{0.45\textwidth}
    \begin{block}{\small Peng et al. 2024 --- GraphRAG Survey (ACM TOIS)}
    \scriptsize
    \textbf{Idée} : première survey systématique formalisant le workflow GraphRAG.

    \vspace{0.3em}
    \textbf{Apports}
    \begin{itemize}\setlength\itemsep{-0.1em}
      \item Formalisation du workflow \textbf{G-Indexing / G-Retrieval / G-Generation}
      \item Panorama des applications (QA, sciences, industrie)
      \item Recensement des benchmarks et des cas industriels
      \item Couverture historique large
    \end{itemize}
    \textbf{Rôle} : ancrage académique --- complément historique du survey Zhang
    \end{block}
  \end{columns}
\end{frame}

\begin{frame}{Papiers lus --- KG juridiques}
  \begin{columns}[T,onlytextwidth]
    \column{0.45\textwidth}
    \begin{block}{\small Belikov \& Raoult 2025 --- KG Cassation FR}
    \scriptsize
    \textbf{Idée} : KG construit à partir de \textbf{2820 pourvois pénaux} via Judilibre.

    \vspace{0.3em}
    \textbf{Apports}
    \begin{itemize}\setlength\itemsep{-0.1em}
      \item Ontologie pénale \textbf{sur-mesure} (LegalRuleML/LKIF jugés insuffisants)
      \item Pattern \textit{"Person A → Crime B → Punishment C → texte/précédent"}
      \item Comparaison \textbf{property graph vs RDF}
      \item Précision 93\% / rappel 89\% (annoncés)
    \end{itemize}
    \textbf{Limites}
    \begin{itemize}\setlength\itemsep{-0.1em}
      \item Métriques peu documentées, pas de code/dataset publié
      \item Limité au pénal --- pas de GraphRAG
    \end{itemize}
    \textbf{Pertinence} : \textbf{seul KG sur la Cassation FR} --- preuve de faisabilité
    \end{block}

    \column{0.45\textwidth}
    \begin{block}{\small D'Amato et al. 2025 --- KG CEDH (violences femmes)}
    \scriptsize
    \textbf{Idée} : comparaison \textbf{bottom-up vs LLM} sur le même corpus (HUDOC, 65 jugements).

    \vspace{0.3em}
    \textbf{Apports}
    \begin{itemize}\setlength\itemsep{-0.1em}
      \item Méthodologie \textbf{CQ-driven} (13 Competency Questions)
      \item Bottom-up : 10\,325 triples, ontologie riche, 0 hallucination
      \item LLM (GPT-4o + Mixtral) : T-box seule, score CQ 61.5\%
      \item Ressource \textbf{FAIR} (Zenodo, GitHub, SPARQL endpoint)
    \end{itemize}
    \textbf{Limites}
    \begin{itemize}\setlength\itemsep{-0.1em}
      \item Anglais uniquement, petit corpus
      \item LLM ne peuple pas d'instances
    \end{itemize}
    \textbf{Pertinence} : méthodologie transposable --- modèle d'évaluation
    \end{block}
  \end{columns}
\end{frame}

% ============================================================
\section{Synthèse}

\begin{frame}{Gaps identifiés}
  \begin{enumerate}
    \item Très peu de \textbf{KG juridiques FR} ouverts et réutilisables
    \item GraphRAG peu évalué sur critères \textbf{juridiquement pertinents} (factualité, traçabilité, revirements)
    \item Ontologies FR fragmentées (ELI, AKN, LKIF\ldots) --- pas unifiées JP + articles
    \item Benchmarks LLMs en droit FR limités (peu de tâches multi-saut)
    \item Peu de travaux sur la \textbf{temporalité} / versioning dans les KG juridiques
  \end{enumerate}

  \vspace{0.5em}
  \keybox{Chaque gap = une porte d'entrée pour la contribution.}
\end{frame}

% ============================================================
\section{Positionnement \& contribution}

\begin{frame}{Positionnement \& contribution}
  \keybox{Construire un KG juridique FR + GraphRAG évalué rigoureusement, avec benchmark des LLMs FR en droit.}

  \vspace{0.3em}
  \textbf{Axes de contribution}
  \begin{enumerate}\footnotesize
    \item \textbf{Benchmark LLMs open-source sur droit FR}
    \begin{itemize}\scriptsize
      \item Gemma 3, Kimi, MiniMax, Llama, Mistral, Qwen\ldots
      \item Zero-shot / few-shot \textbf{et} fine-tuning (LoRA/QLoRA) $\rightarrow$ mesurer le \emph{delta}
    \end{itemize}

    \item \textbf{Ontologie juridique adaptée} (base ELI/AKN/LKIF + temporalité)

    \item \textbf{Construction du KG} : pipeline d'extraction, architecture hybride 3 niveaux, mises à jour incrémentales

    \item \textbf{Évaluation comparative des méthodes de retrieval}
    \begin{itemize}\scriptsize
      \item similarité / logique / GNN / LLM / multi-round / hybride
    \end{itemize}

    \item \textbf{Benchmark de pertinence vs systèmes existants}
    \begin{itemize}\scriptsize
      \item Baselines : RAG naïf, Doctrine.fr, Dalloz IA, LLM seul
      \item Métriques : factualité, traçabilité, recall des fondements, gestion des revirements
    \end{itemize}
  \end{enumerate}

  \vspace{0.3em}
  \footnotesize\textit{Pipeline cohérent : chaque brique publiable indépendamment.}
\end{frame}

% ============================================================
\section{Roadmap}

\begin{frame}{Roadmap --- vue d'ensemble}
  \centering
  \begin{tikzpicture}[scale=0.95, transform shape,
      phase/.style={rectangle, rounded corners, draw=accent, thick, minimum width=1.6cm, minimum height=0.7cm, align=center, font=\scriptsize},
      current/.style={phase, fill=warn!30, draw=warn, very thick},
      done/.style={phase, fill=accentsoft},
      todo/.style={phase}]
    \node[done] (p0) at (0,0) {P0\\État de l'art};
    \node[current] (p1) at (2,0) {P1\\Ontologie\\données};
    \node[todo] (p2) at (4,0) {P2\\Construction\\KG};
    \node[todo] (p3) at (6,0) {P3\\Benchmark\\LLMs};
    \node[todo] (p4) at (8,0) {P4\\Data\\science};
    \node[todo] (p5) at (10,0) {P5\\GraphRAG\\\& retrieval};
    \node[todo] (p6) at (12,0) {P6\\Rédaction};
    \foreach \x/\y in {p0/p1, p1/p2, p2/p3, p3/p4, p4/p5, p5/p6} \draw[-{Stealth}, thick, accent] (\x) -- (\y);
    \node[font=\tiny, warn] at (2, -0.8) {\textbf{ici}};
  \end{tikzpicture}

  \vspace{0.6em}
  \scriptsize
  \begin{tabular}{l l l}
    \toprule
    \textbf{Phase} & \textbf{Objet} & \textbf{Lien contribution} \\
    \midrule
    P0 & Cartographie, gaps & clôture \\
    P1 & Ontologie + Judilibre/Légifrance & Axe 2 \\
    P2 & Extraction LLM, hybride 3 niveaux & Axe 3 \\
    P3 & Benchmark LLMs (zero-shot + fine-tuning) & Axe 1 \\
    P4 & Relations, centralité, communautés & Axe 3 enrichi \\
    P5 & Méthodes de retrieval + benchmark final & Axes 4 + 5 \\
    P6 & Rédaction mémoire & --- \\
    \bottomrule
  \end{tabular}

  \vspace{0.3em}
  \footnotesize\textit{P1 et P3 peuvent avancer en parallèle.}
\end{frame}

\begin{frame}{Prochaines étapes --- 4 à 6 semaines}
  \textbf{Court terme}
  \begin{itemize}
    \item Clôturer P0 : finir fiches GraphRAG + KG juridiques
    \item Poursuivre le fichage GraphRAG (LightRAG, HippoRAG, KAG, StructRAG)
    \item Décider \textbf{périmètre du KG} : JP seule ou JP + articles ?
    \item Explorer Judilibre / Légifrance --- récupérer un échantillon
    \item Choisir base ontologique (ELI + extension ? AKN ?)
  \end{itemize}

  \begin{alertblock}{Points d'arbitrage avec le superviseur}
    \begin{itemize}
      \item Périmètre du KG
      \item Priorité entre les 5 axes de contribution
      \item Accès au cluster GPU (nécessaire pour P3)
    \end{itemize}
  \end{alertblock}
\end{frame}

% ============================================================
\section*{Conclusion}

\begin{frame}{En résumé}
  \begin{itemize}
    \item \textbf{État de l'art quasi bouclé} --- cartographie 50 sources, papiers pivots fichés
    \item \textbf{Position claire} --- pipeline cohérent en 5 axes de contribution
    \item \textbf{Spécificité juridique assumée} --- traçabilité, temporalité, explicabilité
    \item \textbf{Roadmap structurée} en 7 phases avec jalons superviseur
  \end{itemize}

  \vfill
  \centering
  \Large\textbf{Questions ?}
\end{frame}

% ============================================================
\appendix
\begin{frame}{Annexes --- disponibles en backup}
  \begin{itemize}
    \item Tableau complet des 50 sources
    \item Schémas détaillés (pipeline Edge 2024, taxonomie Zhang 2025)
    \item Exemples bruts Judilibre / Légifrance
    \item Modèle temporel envisagé (bi-temporel : \textit{valid time} / \textit{transaction time})
    \item Liens Obsidian / Git
    \item Glossaire juridique
  \end{itemize}
\end{frame}

\end{document}
